\section{The retained mass enters the actual prime-boundary receiver}
\label{sec:mass-theta-prime}

\paragraph{MRT1. The two source moments and their prime class.}
The preceding construction has an arithmetic receiver, which we
now calculate rather than identifying the two source spaces by
analogy. Use the even Schwartz space \(H=\mathcal S(\mathbb R)^{\rm even}\),
with the two original moments
\[
 m_H(h)=\left(h(0),\int_{\mathbb R}h(y)\,dy\right),\quad
 H_0=\ker m_H,\quad R=\mathbb C[\mathbb Q_{>0}^{\times}],\quad
 I=\ker(\operatorname{aug}:R\to\mathbb C).                 \tag{MRT1}
\]
The group basis elements obey \(t_at_b=t_{ab}\), without changing
the rational dilation parameter. The actual moment-zero source
and its algebraic coinvariants are
\[
 M_0=\ker(\operatorname{aug}\otimes m_H),\qquad
 Q_0=M_0/IM_0,\qquad
 \mathcal B=\bigoplus_{p\ {\rm prime}}\mathbb C\ell_p\otimes\mathbb C^2.
                                                               \tag{MRT2}
\]
These are the original periodization source and prime-boundary
coordinates of the published
\href{https://github.com/KokunoYumeto/zeta-function-research-reader/blob/ec73ac350d0cda3d95c0bd1361bede33166b4d34/workbenches/splitzero-tandem/continuations/20260923-original-zeta-return/NOTE.tex}{original-zeta return, FR15--25}.
The periodization and its measure come from
\href{https://arxiv.org/abs/math/9811068v1}{Alain Connes,
\emph{Trace formula in noncommutative geometry and the zeros of the
Riemann zeta function}}, Section III (6)--(19). The following
explicit coefficient map and its mass-completion extension are
derived here; they are not attributed to that article.

Here is the full algebra needed for this receiver. The functions
\[
 h_0(y)=(1-2\pi y^2)e^{-\pi y^2},\qquad
 h_1(y)=2\pi y^2e^{-\pi y^2}
\]
have moments \((1,0)\) and \((0,1)\), respectively. Gaussian
integration and evaluation prove these values. They give
\(H=H_0\oplus\operatorname{span}(h_0,h_1)\), so
\[
 Q_0\cong H_0\oplus(I/I^2)\otimes\mathbb C^2
       \cong H_0\oplus\mathcal B.                         \tag{MRT3}
\]
Indeed \(M_0=(R\otimes H_0)\oplus(I\otimes\mathbb C^2)\);
its multiplication by \(I\) has components
\(I\otimes H_0\) and \(I^2\otimes\mathbb C^2\).
In \(I/I^2\), the identity
\(t_{ab}-1=(t_a-1)+(t_b-1)+(t_a-1)(t_b-1)\), followed by
the prime factorization of a positive rational, expresses every
class as \(\sum_p v_p(a)[t_p-1]\). The linear functionals
\(t_a\mapsto v_p(a)\) kill \(I^2\) and separate these
generators. This proves both isomorphisms and their inverses.

Explicitly the coordinates of \(c=[\sum_a t_a\otimes h_a]\) are
\[
 \Phi c=\sum_a h_a,\qquad
 \beta(c)=\sum_p\ell_p\otimes\sum_a v_p(a)m_H(h_a).          \tag{MRT4}
\]
The defining source condition gives \(\Phi c\in H_0\).
An inverse to MRT3 sends \((h,\sum_p\ell_p\otimes u_p)\)
to
\[
 [1\otimes h]+\sum_p[(t_p-1)\otimes
                    (u_{p,0}h_0+u_{p,1}h_1)] .             \tag{MRT5}
\]
There are only finitely many nonzero \(u_p\). Changing the
moment lift changes this class by \(IM_0\). These calculations
verify all quotient relations; the receiver is not merely a
map on arbitrarily selected representatives.

\paragraph{MRT2. The heat coefficient source has an exact two-moment map.}
Take an even finitely supported coefficient sequence \(a\), with
\(a_{-n}=a_n\), no lower coefficients for this specified arrow,
and the unchanged \(T>0\). Its heat function \(h=A_Ta\) of
MCB12 belongs to \(H\). Direct evaluation and integration give
\[
 q_T(a)=\sum_n a_n k_T(n),\qquad m(a)=\sum_n a_n,\qquad
 m_H(A_Ta)=(q_T(a),m(a)).                                  \tag{MRT6}
\]
Consequently, for every original arithmetic prime \(p\), the map
\[
 B_{p,T}(a)=[(t_p-1)\otimes A_Ta]\in Q_0
\]
has the exact coordinates
\[
 (\Phi,\beta)B_{p,T}(a)
        =(0,\ell_p\otimes(q_T(a),m(a))).                   \tag{MRT7}
\]
The first zero follows by subtraction of the same function, and
the second follows from \(v_p(p)=1\), \(v_p(1)=0\).
This is the connecting map from the actual integer heat source
to the existing prime-boundary source.

Write \(k_0=k_T(0)\), \(k_1=k_T(1)\); these symbols retain
the values \((4\pi T)^{-1/2}\) and
\((4\pi T)^{-1/2}e^{-1/(4T)}\). For any pair \((u_0,u_1)\),
put
\[
 a_0=\frac{u_0-k_1u_1}{k_0-k_1},\qquad
 a_{1}=a_{-1}=\frac{k_0u_1-u_0}{2(k_0-k_1)},\qquad
 a_n=0\quad(|n|>1).                                      \tag{MRT8}
\]
Since \(k_0>k_1\), this is defined. Substitution gives
\(q_T(a)=u_0\) and \(m(a)=u_1\). Thus this actual finite
source attains both prime-boundary moments, with their original
coefficients and Gaussian constants. It is not confined to a
single prescribed moment ratio.

\paragraph{MRT3. Completion and supported-zero action commute with this map.}
Let \(\widehat V_E^{\rm even,top}\) be the closed subspace
\((a,0,m)\) of MCB5 with even \(a\). The samples
\((k_T(n))_n\) lie in \(\ell^2\); hence Cauchy--Schwarz gives
\[
 |q_T(a)|\le
 \left(\sum_n k_T(n)^2\right)^{1/2}\|a\|_2.              \tag{MRT9}
\]
The same formula with an independent last coordinate extends MRT7:
\[
 \widehat B_{p,T}(a,0,m)=
 b\bigl(\ell_p\otimes(q_T(a),m)\bigr),                   \tag{MRT10}
\]
where \(b\) is the explicit map MRT5 on its second summand.
The two scalar coordinates are continuous in the literal graph
norm. No new norm is assigned to the adelic quotient in this
statement. Even finite graph sources are dense in this subspace:
in MCB6 choose the correcting sets in opposite pairs, retaining
the same total correction and its vanishing squared norm.

The original supported-zero operator now has the following
arithmetic receiving matrix, with all constants retained:
\[
 J_T=\begin{pmatrix}0&k_T(0)\\0&1\end{pmatrix},\qquad
 J_T^2=J_T,\qquad
 \widehat B_{p,T}\widehat E
   =b\,(\operatorname{id}_{\mathbb C\ell_p}\otimes J_T)
                         \beta\widehat B_{p,T}.            \tag{MRT11}
\]
Indeed \(\widehat E(a,0,m)=(m\delta_0,0,m)\), whose two
moments are \((k_T(0)m,m)\). This proves MRT11 exactly;
it does not assert a not-yet-defined multiplication on all of
\(Q_0\). On the completed pure-mass vector \((0,0,1)\),
MRT10 gives \(b(\ell_p\otimes(0,1))\); after \(\widehat E\)
it gives \(b(\ell_p\otimes(k_T(0),1))\). The boundary action
is consequently visible in the actual two-moment coordinates.

\paragraph{MRT4. Relation to the original periodization seminorm.}
Keep exactly the original periodization and its measure:
\[
 \mathcal E(c)(x)=x^{1/2}\,2\sum_{n\ge1}(\Phi c)(nx),\qquad
 \|c\|_\delta^2=\int_0^\infty|\mathcal E(c)(x)|^2
                 (1+\log^2x)^{\delta/2}\frac{dx}{x},
 \quad\delta>1.                                          \tag{MRT12}
\]
Connes's compact-unit measure has mass one in this formula; with
mass \(c_K\) the whole displayed integral has the factor \(c_K\),
which is retained. For every vector in MRT10, \(\Phi=0\),
so \(\mathcal E=0\) and this exact seminorm is zero. Whenever
\((q_T(a),m)\ne(0,0)\), MRT3 proves that its boundary class is
nonzero. This verifies on the actual Gaussian coefficient source
the failure of its two moments to descend through the Hausdorff
quotient of MRT12. The graph completion MCB5 repairs the domain
of supported-zero multiplication; it does not turn MRT12 into
a positive norm on these prime classes. The maps MRT7--12 retain
both facts and give their exact connection.

\paragraph{MRT5. The theta family is an admitted source, with two evaluated moments.}
Let \(x>0\), independently of the heat time \(T\), and put
\[
 a^{x}_n=e^{-\pi n^2x},\qquad
 \vartheta(x)=\sum_{n\in\mathbb Z}e^{-\pi n^2x},\qquad
 d_T=\frac1{4\pi T},\qquad k_T(0)=\sqrt{d_T}.             \tag{MRT13}
\]
These original weights are even and absolutely summable. Their
finite truncations converge in the graph norm to
\((a^x,0,\vartheta(x))\): both the squared tails and the mass
tails tend to zero by Gaussian convergence. Unlike the unweighted
integer comb, this is an admitted vector of MCB5.

Moreover \(A_Ta^x\) is an even Schwartz function. To see every
Schwartz bound directly, each derivative of \(k_T(y-n)\) is a
fixed polynomial in \(y-n\) times its Gaussian. For any integer
\(M\ge0\), use
\(1+|y|\le(1+|n|)(1+|y-n|)\). The supremum of the resulting
polynomial times that Gaussian is finite, and its sum against
\((1+|n|)^M e^{-\pi n^2x}\) is finite. Uniform derivative
convergence proves the assertion. Hence MRT7 applies directly
to this actual Schwartz function, not just to its completed image.
Evaluation and integration in MRT6 yield
\[
 \beta B_{p,T}(a^x)=
 \ell_p\otimes\left(
    \sqrt{d_T}\,\vartheta(x+d_T),\ \vartheta(x)\right).
                                                               \tag{MRT14}
\]
The product of the two Gaussian factors proves the first entry,
and integration of each translate proves the second. After the
supported-zero action the same class has moments
\[
 \ell_p\otimes\left(
     \sqrt{d_T}\,\vartheta(x),\ \vartheta(x)\right).       \tag{MRT15}
\]
Thus the change in this prime-boundary coordinate, in the order
before minus after, is exactly
\[
 \ell_p\otimes\left(
   \sqrt{d_T}[\vartheta(x+d_T)-\vartheta(x)],\ 0\right).
                                                               \tag{MRT16}
\]
It is nonzero for every \(x,T>0\), since every pair of nonzero
integer Gaussian terms strictly decreases when its parameter
increases by \(d_T>0\). This evaluates an actual arithmetic
consequence of the supported-zero operation in the retained source.

\paragraph{MRT6. Return to the original zeta function, with the marked term kept.}
The single integer-zero source is \(\delta_0\), with coefficient
one, mass one and heat moments \((\sqrt{d_T},1)\). Keep it
as a separate summand. The remaining pair from MRT14 is
\[
 \left(\sqrt{d_T}[\vartheta(x+d_T)-1],\ \vartheta(x)-1\right).
                                                               \tag{MRT17}
\]
For \(\Re s>1\), its second Mellin integral is exactly
\[
 \int_0^\infty[\vartheta(x)-1]x^{s/2-1}dx
   =2\pi^{-s/2}\Gamma(s/2)\zeta(s),\qquad
 \zeta(s)=\frac{\pi^{s/2}}{2\Gamma(s/2)}
                \int_0^\infty[\vartheta(x)-1]x^{s/2-1}dx.
                                                               \tag{MRT18}
\]
Indeed the paired integers \(\pm n\) give the factor two;
the substitution \(y=\pi n^2x\) gives
\(\pi^{-s/2}n^{-s}\Gamma(s/2)\); absolute interchange follows
from \(\sum n^{-\Re s}<\infty\). The marked constant has
no Mellin convergence half-plane on \((0,\infty)\) and remains
separate, rather than being assigned zero. Poisson reflection and
splitting at one retain two distinct endpoint contributions: the
zero source coefficient gives \(-1/s\), and the zero Fourier
coefficient gives \(1/(s-1)\), as derived in LH15. Thus original \(\zeta\)
and all its continuation factors remain explicit.

The first coordinate also has an evaluated Mellin transform:
for \(\Re s>0\),
\[
 \int_0^\infty\sqrt{d_T}[\vartheta(x+d_T)-1]x^{s/2-1}dx
  =2\sqrt{d_T}\,\pi^{-s/2}\Gamma(s/2)
         \sum_{n\ge1}e^{-\pi d_Tn^2}n^{-s}.                \tag{MRT19}
\]
Here the same calculation is absolutely convergent already for
\(\Re s>0\), because the Gaussian coefficient controls the
integer sum. The full Dirichlet series in MRT19 is entire in
\(s\), by uniform Gaussian convergence on compact sets; the
Gamma factor is retained. In the common half-plane \(\Re s>1\),
the Mellin transform of MRT16 is therefore
\[
 2\sqrt{d_T}\,\pi^{-s/2}\Gamma(s/2)
       \left[\sum_{n\ge1}e^{-\pi d_Tn^2}n^{-s}-\zeta(s)\right]
       \ell_p\otimes(1,0).                               \tag{MRT20}
\]
Every term in this difference has now been identified. In particular
the zero-seminorm prime-boundary source contains a family whose
retained mass returns the original zeta function. This is a proved
receiver and a proved loss of specified data under MRT12; it is
not an evaluation of the full Weil pairing on a newly invented norm.
